%% tables/tab_elements.tex

\begin{table*}
\centering
\caption{
  Osculating elements (heliocentric ECLIPJ2000, epoch J2000.0)
  of benchmark asteroids used in this work.
  \label{tab:elements}
}
\begin{tabular}{lccccccc}
\toprule
Body & $a$ (AU) & $e$ & $i$ (deg) & $\Omega$ (deg) & $\omega$ (deg) & $M$ (deg) \\
\midrule
(1) Ceres        & 2.7665 & 0.0784 & 10.583 &  80.494 &  73.923 &   6.177 \\
(99942) Apophis  & 0.9223 & 0.1914 &  3.331 & 204.657 & 126.065 & 232.181 \\
(3200) Phaethon  & 1.2715 & 0.8902 & 22.129 & 265.510 & 321.894 & 142.572 \\
(309704) Baruffetti & 3.0586 & 0.0503 & 10.629 &  52.331 &  26.303 & 265.870 \\
\bottomrule
\end{tabular}

\medskip
\parbox{0.95\textwidth}{\footnotesize
Source: JPL Horizons (retrieved 2025 March, epoch J2000.0).
Orbital elements are computed from the heliocentric ICRF state vectors
by rotating to ECLIPJ2000 (obliquity $\varepsilon=23.4393^\circ$).
$\mu = \mathrm{GMS} = k^2 = 2.9591\times10^{-4}$\,AU$^3$\,day$^{-2}$.
No planetary perturbations are applied in the two-body benchmarks
(Section~\ref{sec:benchmark}).
}
\end{table*}
